\section{Implementation}
\label{sec:impl}

\subsection{Execution}

CMCA runs on Harbor~\citep{harbor}, pinned at version 0.24.0.
Each harness runs as a Harbor installed agent that installs the pinned harness version, copies a prebuilt injector into the container, writes the resolved engine configuration and collects the injector log after the run.
Profile, seed, parity set and attempt index are passed as agent arguments and never appear in the task instruction.
Trajectories are stored in Harbor's ATIF format.
Every injector writes one JSON record per eligible call (decision, fault type, payload, capability class, argument hash, occurrence), and every Harbor job writes a manifest with the harness and model versions, injector build hash, catalog version, profile hash, seed and task image digests.

A shared TypeScript engine implements the schedule; the MCP proxy and every native injector call it, so the same call key yields the same decision on every layer.
Property tests check that decisions are deterministic, that the observed rate over 100,000 random keys stays within 0.5 percentage points of the configured rate, and that caps hold under concurrent calls.

\subsection{Harness integrations}

\begin{itemize}
\item \textbf{Claude Code}: a plugin whose handlers act before and after each built-in tool call.
\item \textbf{Codex CLI}: command hooks (\code{PreToolUse}, \code{PostToolUse}). Codex speaks only the Responses API, so a bridge we wrote, running in the container, translates its requests to the Messages API. Codex itself is unmodified.
\item \textbf{Pi}: a TypeScript extension on its tool-call events.
\item \textbf{OpenCode}: a plugin on its tool execution hooks.
\item \textbf{fx}: service track only; the proxy injects, because fx~0.0.13 exposes no interception point for its built-in tools.
\item \textbf{fxh}: native track only; a fork of fx~0.0.13 maintained by the author that adds pre- and post-execution hooks for built-in tools, loaded with the same injector contract as the others.
\end{itemize}

\subsection{Conformance}
\label{sec:conformance}

No board is published for a harness version that has not passed conformance.
For each harness version, fault type and capability class, a probe task makes one targeted call under a profile that forces the fault on its first eligible occurrence.
The probe checks that the injector logged the decision, that the payload reached the model's observation, whether the harness added attribution text, and that side effects match the phase (none before execution, present after).
Each probe runs three times, and a fault type enters a harness's parity set only if all runs pass.
The output is a machine-readable capability file per harness version, from which parity sets are generated, never written by hand.

Conformance and the delivery audit found integration behavior that changes what a model sees.
Claude Code's mapper renders a fixed success line for writes and edits, so an empty result cannot reach the model on those classes; they are declared parity gaps.
In Claude Code and Pi, the first injectors returned every failed shell command's result unchanged, so post-execution faults on commands that exited non-zero were never delivered; a probe for failed commands now guards this path.
Claude Code saves large command outputs to a file before its post-execution hook runs, so the first injector faulted only an empty preview; the injector was fixed during the run (Table~\ref{tab:deviations}).
Codex fires its post-execution hook for a command still running when the tool returns only when the process ends, on a later poll, so a fault on such a command can miss the model.
Each finding is recorded with its evidence in the repository's verification log.
